% =====================================================================
% Введение (по стилю HW_10_PRESENTATION_NP.tex)
% =====================================================================
\begin{frame}{Введение}

\textbf{Актуальность:} солисты-пианисты при подготовке к концерту
нуждаются в репетициях с оркестровым сопровождением, однако живой
оркестр недоступен на каждой репетиции. Фиксированная фонограмма не
подстраивается под темп и \emph{rubato} солиста.

\bigskip

\textbf{Проблема:} построить виртуальный оркестр, который в реальном
времени отслеживает позицию исполнителя в партитуре и упреждающе
управляет темпом аккомпанемента (бюджет задержки $\le 20$\,мс ---
порог восприятия рассинхронизации в ансамбле, Cont 2010).

\bigskip

\textbf{Подход:} классический трекер партитуры (HMM\,$+$\,OLTW)
$+$ \emph{лёгкий RL-агент}, который видит состояние трекера и
управляет темпом на горизонте 200--500\,мс.

\bigskip

\textbf{Цель доклада:} разобрать математическую часть baseline,
обосновать формулу награды RL-агента и положение работы относительно
литературы.

\end{frame}
